%% ma: L10-B14-0015
%% lop: 10
%% chuong: 5
%% bai: 14
%% dang: 10.14.2
%% loai: TL
%% muc: TH
%% dap_an: "-"
%% nguon: "chuyen-de-cac-so-dac-trung-cua-mau-so-lieu-khong-ghep-nhom-toan-10-chuong-trinh-moi.pdf (GV. Phan Nhật Linh) (P04), Bài 03 Dạng 1 Phần I Câu 5"
%% kien_thuc: "Lớp 10 Bài 13 (tứ phân vị), Bài 14 (khoảng tứ phân vị)"
%% trang_thai: chinh-thuc
%% ghi_chu: "Câu nằm ở Phần I nhưng đề gốc không có phương án, đã nhập dạng tự luận. Tự duyệt theo yêu cầu giáo viên 22/10/2026."
\begin{ex}
Chiều cao (đơn vị: cm) của các bạn trong tổ của Nam như sau:
\begin{center}
\begin{tabular}{|c|c|c|c|c|c|c|c|c|c|}\hline
$165$ & $168$ & $157$ & $162$ & $165$ & $165$ & $179$ & $148$ & $170$ & $167$\\\hline
\end{tabular}
\end{center}
Tính khoảng tứ phân vị của mẫu số liệu. Khoảng tứ phân vị có bị ảnh hưởng bởi chiều cao của bạn cao nhất, bạn thấp nhất tổ không?
\loigiai{Sắp xếp: $148$; $157$; $162$; $165$; $165$; $165$; $167$; $168$; $170$; $179$.

$Q_2=\dfrac{165+165}2=165$; nửa dưới $148$; $157$; $162$; $165$; $165$ nên $Q_1=162$; nửa trên $165$; $167$; $168$; $170$; $179$ nên $Q_3=168$.

$\Delta_Q=168-162=6$.

Khoảng tứ phân vị chỉ phụ thuộc vào $50\%$ số liệu ở giữa nên không bị ảnh hưởng bởi giá trị lớn nhất và nhỏ nhất.}
\end{ex}
